\clearpage
\ssscp{Apicals *d, *z, *R, *l, *j}{07-04-01-02}
Another feature which distinguishes \tsc{Ambon-Seram} from neighbouring
subgroups is the behaviour of the apicals *d,
*z, *R, *l, and *j.
\tsc{Ambon-Seram} gives the strong impression of merging *R/*l/*j
> **l, while they are kept distinct in \tsc{Sula-Buru}
and \tsc{Seram-Tanimbar-Bomberai} (with the exception of some
lower branches of \tsc{Seram-Tanimbar-Bomberai}).
There is some phonological conditioning in \tsc{Ambon-Seram}
languages, such that the reflexes of *l
and *R contiguous to [+LOW] vowel \it{a},
may be different than contiguous to [+HIGH] vowels \it{i, u}.
See \srf{10-01} and \citet[23--33]{Stresemann1927} for detailed
discussion of the different patterns observed in the data.
Yet like the behaviour of PMP *j,
different conditioned reflexes and lexically specific exceptions
can also be seen,
meaning *R and *l were not fully
and completely merged in \tsc{Ambon-Seram}.
The reflexes of the apicals *d/*z/*R/*l/*j in subgroups of central
Maluku and two nodes of \tsc{Ambon-Seram} are summarised in \trf{07-09}.

%\dirtree{%%The percent after { is needed
%.1 {\tsc{Ambon-Seram}}.
		%.2 {Boano}.
		%.2 {\tsc{Nuclear Ambon-Seram}}.
				%.3 {\tsc{Ambon-West Seram}}.
				%.3 {\tsc{Pataki-Manusela}}.
		%.2 {\tsc{Seti-Benggoi}}.
%}%Ensure there is no carriage return before this }

\begin{table}\tcp{Reflexes and mergers of apicals in central Maluku}{07-09}
	\begin{threeparttable}\st{6pt}
	\begin{tabular}{lllllllll}\lsptoprule
PMP &{\hr}*d&{\hr}*z&{\hr}*R&{\hr}*l&{\hr}*j\\ \midrule
\tsc{Sula-Buru}&**d&**z&**h&**l{\cdr}&\mc{2}{l}{**y, **l, {\0}}\\ \midrule
%\hp{---}Boano&\hp{**}l&\hp{**}l&\hp{**}n&\hp{**}l&\hp{**}n\\
%\hp{---}\tsc{Nuclear Ambon-Seram}&**r&**r&**l&**l&**l\\
%\hp{{---}{---}}\tsc{Ambon-West Seram}&**l&**l&**l&**l&**l\\
%\hp{{---}{---}}\tsc{Patakai-Manusela}&**h&**h&**l&**l&**l\\
%\hp{---}\tsc{Seti-Benggoi}&\hp{**}l&\hp{**}l&\hp{**}l, h&\hp{**}l&\hp{**}l\\ \midrule
\multirow{6}{*}{\includegraphics{Tab07-09.pdf}}
&**r{\cor}&**r{\cor}&**R&**l{\cdr}&**l{\cdr}&\\
&\hp{**}l{\cdr}&\hp{**}l{\cdr}&\hp{**}n{\cye}&\hp{**}l{\cdr}&\hp{**}n{\cye}&\\
&**r{\cor}&**r{\cor}&**l{\cdr}&**l{\cdr}&**l{\cdr}&\\
&**l{\cdr}&**l{\cdr}&**l{\cdr}&**l{\cdr}&**l{\cdr}&\\
&**h{\crd}&**h{\crd}&**l{\cdr}&**l{\cdr}&**l{\cdr}&\\
&\hp{**}l{\cdr}&\hp{**}l{\cdr}&\hp{**}l, h&\hp{**}l{\cdr}&\hp{**}l{\cdr}&\\ \midrule
\tsc{Seram-Tanimbar-Bomberai}
&**d{\cbl}&**d{\cbl}&**r{\cor}&**l{\cdr}&**j&\\ \midrule
Banda&\hp{**}r{\cor}&\hp{**}r{\cor}&\hp{**}r{\cor}&\hp{**}l{\cdr}&\hp{**}l{\cdr}&\\
		\lspbottomrule
	\end{tabular}
		\begin{tablenotes}\tnsize
			\item[†] {While \tsc{Seram-Tanimbar-Bomberai} does have initial *d/*z > **d,
								the data are a little ambiguous for medial *d/*z > **d.
								See \srf{11-01-01} for discussion.}
		\end{tablenotes}
	\end{threeparttable}
\end{table}

\trf{07-10} shows the reflexes of *l in central Maluku languages.
(For more data see also *bulan in \trf{07-01},
*qaləjaw ‘day’ in \trf{07-05},
*layaR ‘sail’ in \trf{07-10},
and *təlu ‘three’ in \trf{07-14}.) The data show that while *l
= \it{l} is quite stable in \tsc{Sula-Buru} (other than a split
in Ambelau to \it{l,
h}) and \tsc{Seram-Tanimbar-Bomberai}, the reflexes of *l in
\tsc{Ambon-Seram} have great diversity in the daughters.
One can also see that the reflexes of *l in \tsc{Ambon-Seram}
do not always align with the reflexes of *R in \trf{07-11}.
Conditioned reflexes of *l are underlined; underlining of a vowel
or vowels indicates a loss of *l.

\trf{07-11} shows reflexes of PMP *R,
with *R > **h in \tsc{Sula-Buru},
and *R > **l in \tsc{Ambon-Seram}.
(For additional reflexes of PMP *R,
see also *timuR ‘east monsoon’ in \trf{07-02}.) As in \trf{07-10},
keep in mind that \tsc{Ambon-Seram} languages for the most part
have merged PMP apicals *d/*z/*R/*l/*j > **l,
and different reflexes in languages like Larike are phonologically
conditioned by the vowels.
There is also assimilation at play in some languages when a word contains multiple liquids.
(Note: +t means the language appears to have retained the gloss
of ‘thorn’, alongside the more widespread gloss of ‘bone’ in this region.)
Conditioned reflexes of *R are underlined.

\trf{07-12} shows the lexically specific split of PMP *j > {\0}
in \tsc{Sula-Buru} and *j > **l in \tsc{Ambon-Seram} in PMP *ŋajan ‘name’
and *huaji ‘younger sibling’.
Reflexes of *ŋijuŋ/*qijuŋ/*ujuŋ ‘nose’ also appear to show *j
> {\0} in \tsc{Sula-Buru} but *j > **l in \tsc{Ambon-Seram}.
However, it is not clear whether all the words in the ‘nose’
column in \trf{07-12} are from the same etymon.
The vowels of the \tsc{Sula-Buru} words are not regular from
any of the reconstructions.
Similarly, there are irregularities in the medial consonants
in the \tsc{Ambon-Seram} words which point to additional complexities
in this cognate set.
Some \tsc{Ambon-Seram} words with medial \it{n} may be explicable
as progressive assimilation whereby **l > \it{n} /nV{\gap}.

 In relation to the complexities in \tsc{Ambon-Seram},
note that the retentions of *j are not the same in these three
words for several of the languages,
making broad generalisations about splits
and mergers in \tsc{Ambon-Seram} difficult.
Conditioned reflexes of *j are underlined.

\begin{table}[p]\centering\tcp{Central Maluku *l}{07-10}
\begin{adjustbox}{totalheight=\textheight}
	\begin{threeparttable} %\st{2pt}
	\begin{tabular}{llllll}\lsptoprule
*gloss	&	sky	&	rat, mouse	&	path, road	&	moon	&	ear\\
PMP	&	*\tbr{l}aŋit	&	*(ba/ma)\tbr{l}abaw	&	*za\tbr{l}an	&	*bu\tbr{l}an	&	*ta\tbr{l}iŋa\\ \midrule
%\tsc{Sula-Buru}\\ \midrule
Mangole	&	{\hr}\tbr{l}aŋ, \tbr{l}aŋi	&	{\hr}faʔafa, fafa	&	\cg{(nafakau)}	&	\cg{(fasina)}	&	{\hr}ta\tbr{l}iŋa\\
Sanana	&	{\hr}\tbr{l}an	&		&	{\hr}yaa\su{Δ}	&	\cg{(fasina)}	&	{\hr}te\tbr{l}ina\\
Lisela	&	{\hr}\tbr{l}aŋit, \tbr{l}aŋiʔ	&	{\hr}b\tbr{l}afun pait\su{†}	&	\cg{(tuhu-n)}	&	{\hr}fu\tbr{l}an	&	{\hr}k\tbr{l}iŋa-n\\
Buru	&	{\hr}\tbr{l}aŋit	&	\cg{(titasit)}	&	\cg{(toho-n)}	&	{\hr}fu\tbr{l}an	&	{\hr}et\tbr{l}iŋa-n\\
Ambelau	&	{\hr}\tbr{l}anir-e	&	\cg{(pue)}	&	{\hr}la\tbr{h}-lea	&	{\hr}hu\tbr{l}a	&	{\hr}e\tbr{h}riya-ni (\it{met.})\\ \midrule
%\tsc{Ambon-Seram}\\ \midrule
Kayeli	&	{\hr}\tbr{l}anit	&	\cg{(egboti)}	&	{\hr}la\tbr{l}an	&	{\hr}bu\tbr{l}an	&	{\hr}t\tbr{l}ila-ni\\
Boano	&	{\hr}\tbr{n}anit-e	&		&	{\hr}la\tbr{l}an-e	&	{\hr}hu\tbr{n}an-e	&	\\
Luhu	&	{\hr}\tbr{l}anit-e	&	\cg{(puene)}	&	{\hr}la\tbr{l}an	&	{\hr}hu\tbr{l}an-e	&	{\hr}te\tbr{l}ina\\
Piru	&	{\hr}\tbr{l}anit	&	\cg{(puen)}	&	{\hr}la\tbr{l}an	&	{\hr}hu\tbr{l}an	&	{\hr}ta\tbr{l}ina\\
Larike	&	{\hr}\tbr{l}anit-a	&	{\hr}un{\hr}\tbr{\und{d}}aha	&	{\hr}la\tbr{l}an-u	&	{\hr}hu\tbr{\und{d}}an-u	&	{\hr}te\tbr{\und{d}}ina\\
Asilulu	&	{\hr}\tbr{l}anit	&	{\hr}ma\tbr{l}aha	&	{\hr}la\tbr{l}an	&	{\hr}hu\tbr{l}an	&	{\hr}ta\tbr{l}ina\\
Paulohi	&	{\hr}\tbr{l}anit-e	&	{\hr}ma\tbr{l}aha	&	{\hr}l\und{aa} kine	&	{\hr}hu\tbr{l}an-e	&	{\hr}te\tbr{\und{ʀ}}ina {\tl} ti\tbr{\und{ʀ}}ina\\
Amahai	&	{\hr}\tbr{l}aɲi-o\tbr{l}o	&	{\hr}mu\tbr{\und{r}}aha	&	{\hr}la\tbr{l}a ira	&		&	{\hr}e\tbr{\und{r}}na\\
Kaibobo	&	{\hr}\tbr{l}anit	&	{\hr}ma\tbr{l}aha	&	{\hr}la\tbr{l}an	&	{\hr}hu\tbr{\und{r}}an	&	{\hr}ta\tbr{\und{r}}ina\\
Kamarian	&	{\hr}\tbr{l}anit	&	{\hr}ma\tbr{l}aha	&	{\hr}la\tbr{l}an	&	{\hr}hu\tbr{\und{r}}an	&	{\hr}ta\tbr{\und{r}}ina\\
Haruku	&	{\hr}\tbr{l}anit-o	&	{\hr}ma\tbr{l}aha	&	{\hr}la\tbr{l}an-o	&	{\hr}hu\tbr{\und{r}}an-o	&	{\hr}ta\tbr{\und{r}}ina\\
Saparua	&	{\hr}\tbr{l}anit-ol	&	{\hr}mo\tbr{l}ahal	&	{\hr}la\tbr{l}an-o	&	{\hr}u\tbr{l}an-o	&	{\hr}ta\tbr{l}ina\\
Nusalaut	&	{\hr}\tbr{l}ani-ol	&	{\hr}mo\tbr{l}ahal	&	{\hr}la\tbr{l}an-o	&	{\hr}hu\tbr{l}an-o	&	{\hr}ta\tbr{\und{r}}ina\\
N. Wemale	&	{\hr}\tbr{l}anit-e	&	{\hr}m\und{aa}pa	&	{\hr}l\und{a} rina (< **laa)	&		&	{\hr}t\und{ii}na\\
Sepa	&	{\hr}\tbr{l}aant-o	&	{\hr}m\tbr{l}aha	&	{\hr}l\und{aa} tina	&	{\hr}h\und{ia}n-o	&	{\hr}ti\tbr{\und{r}}ina\\
Tamilou	&	{\hr}\tbr{l}aant-o	&	{\hr}ma\tbr{\und{y}}aha	&	{\hr}l\und{aa} tina	&	{\hr}h\und{ia}n-o	&	{\hr}ti\tbr{\und{n}}na\\
Tehoru	&	{\hr}\tbr{l}aant-o	&	{\hr}ma\tbr{\und{y}}aha	&	{\hr}l\und{aa} tina	&	{\hr}h\und{ia}n-o	&	\\
Lisabata	&		&	{\hr}ma\tbr{l}aha	&		&	{\hr}hu\tbr{l}an-e	&	\\
Naka{\Q}ela	&		&	{\hr}ma\tbr{l}ahaw-e	&		&	{\hr}hu\tbr{l}an-e	&	\\
Watui	&		&	{\hr}m\und{a}boʔe	&		&	{\hr}pu\tbr{l}an-e	&	\\
N. Alune	&	{\hr}\tbr{l}anit-e	&	{\hr}ma\tbr{l}abakw-e	&	{\hr}la\tbr{l}an-e	&	{\hr}bu\tbr{l}an-e	&	{\hr}t\tbr{l}ina\\
W. Alune	&	{\hr}\tbr{l}anit	&	{\hr}ma\tbr{\und{r}}awak-a	&	{\hr}la\tbr{\und{r}}an	&	{\hr}bu\tbr{\und{r}}an	&	{\hr}ta\tbr{l}ina\\
Yalahatan	&	{\hr}\tbr{l}ant-e	&	{\hr}ma\tbr{l}aha{\tl}ma\tbr{l}afa	&	{\hr}ra\tbr{\und{r}}an-e	&	{\hr}hu\tbr{l}an-e	&	{\hr}ti\tbr{\und{r}}ina-i\\
S. Nuaulu	&	{\hr}\tbr{n}ant-e	&	{\hr}m\tbr{n}aha	&	\cg{(arena)}	&	{\hr}hu\tbr{n}an-e	&	{\hr}tina-i\\
Manusela	&	{\hr}\tbr{l}ali-a	&	{\hr}ma\tbr{\und{i}}aha	&	{\hr}ha\tbr{l}a	&	\cg{(umala)}	&	{\hr}tina-nia\\
Liana-Seti	&	\cg{(lea)}	&		&	{\hr}la\tbr{l}am	&	\cg{(umala)}	&	{\hr}tina\\  \midrule
Banda	&	{\hr}\tbr{l}aŋit	&		&	\cg{(eran)}	&	{\hr}wu\tbr{l}an	&	\cg{(fulun, tilu-)}\\ \midrule
%\tsc{Seram-Tanimbar-Bomberai}\\  \midrule
Bobot	&	{\hr}\tbr{l}akit	&		&	{\hr}lo\tbr{l}an	&	{\hr}vu\tbr{l}an	&	{\hr}ta\tbr{l}ikan\\
Geser	&	{\hr}\tbr{l}aŋit	&	\cg{(kalufa)}	&	{\hr}la\tbr{l}an	&	{\hr}u\tbr{l}an	&	{\hr}ti\tbr{l}iŋa\\
Gorom	&	{\hr}\tbr{l}aŋit	&	\cg{(aruha)}	&	{\hr}la\tbr{l}an	&	{\hr}u\tbr{l}an	&	{\hr}te\tbr{l}iŋa\\
Watubela	&	{\hr}\tbr{l}aŋit	&		&		&	{\hr}u\tbr{l}an	&	{\hr}te\tbr{l}ŋa-\\
Kowiai	&	{\hr}\tbr{r}aŋgit	&	\cg{(asaɸai)}	&	{\hr}la\tbr{r}aʔai	&	{\hr}ɸu\tbr{r}an	&	{\hr}te\tbr{r}iga\\
Yamdena	&	{\hr}\tbr{l}aŋit	&	\cg{(kawat)}	&	{\hr}da\tbr{l}nen	&	{\hr}bu\tbr{l}an	&	{\hr}t\tbr{l}iŋa-ŋw\\
Kei	&	{\hr}\tbr{l}anit	&	\cg{(keru)}	&	\cg{(laŋay)}	&	{\hr}vuan	&	\cg{(aru-n)}\\
Uruangnirin	&	{\hr}\tbr{n}agit	&	\cg{(yatuat)}	&	\cg{(banamate)}	&	{\hr}pu\tbr{n}aŋ	&	{\hr}ta\tbr{n}iŋa-n\\
Sekar	&	\cg{(laŋit <Mly.)}	&	\cg{(kufiata)}	&	{\hr}na\tbr{n}am	&	{\hr}bu\tbr{n}an	&	{\hr}ta\tbr{n}iga-n\\
Arguni	&	{\hr}\tbr{r}ágit	&	\cg{(adumbár)}	&	{\hr}ra\tbr{r}ín	&	{\hr}pú\tbr{r}in	&	{\hr}tigyá\tbr{n}áne\\
Irarutu	&	{\hr}\tbr{r}agt	&	\cg{(sfe)}	&	{\hr}ra\tbr{r}n	&	\cg{(seba)}	&	{\hr}tg\tbr{r}a\\
		\lspbottomrule
	\end{tabular}
		\begin{tablenotes}\tnsize
			\item[†] {Lisela \it{blafun pait} ‘rat, mouse’ is irregular or a loan.
								Unattested *[flafa] would be the regular reflex.}
		\end{tablenotes}
	\end{threeparttable}
\end{adjustbox}
\end{table}

\begin{table}[p]\centering\tcp{Central Maluku *R}{07-11}
\begin{adjustbox}{totalheight=\textheight}
	\begin{threeparttable}
	\st{3.5pt}\begin{tabular}{lllllll}\lsptoprule
local gl.	&	west monsoon	&	blood	&	house	&	bone, thorn	&	new, young	&	sail\\
PMP	&	*haba\tbr{R}at	&	*da\tbr{R}aq	&	*\tbr{R}umaq	&	*du\tbr{R}i	&	*baqə\tbr{R}u	&	*laya\tbr{R}\\ \midrule
%\tsc{Sula-Buru} *R > \it{h}\\ \midrule
Mangole	&		&	\cg{(poo)}	&	{\hr}uma	&	{\hr}loi \it{(+t)}	&	{\hr}feu	&	\\
Sanana	&	\cg{(bara <Mly)}	&	\cg{(poa)}	&	{\hr}uma	&	{\hr}hoi \it{(+t)}	&	{\hr}feu	&	\cg{(soba)}\\
Lisela	&	{\hr}ful fa\tbr{h}at	&	{\hr}ra\tbr{h}a-n	&	{\hr}\tbr{h}uma	&	{\hr}ro\tbr{h}i-n	&	{\hr}fe\tbr{h}u-t	&	{\hr}raa-n \cg{(irr.)}\\
Buru	&	{\hr}ful fa\tbr{h}at	&	{\hr}ra\tbr{h}a-n	&	{\hr}\tbr{h}uma	&	{\hr}ro\tbr{h}i-n	&	{\hr}fe\tbr{h}u, fe\tbr{h}u-t	&	{\hr}laa\\
Ambelau	&	{\hr}ba\tbr{h}ar-e	&	{\hr}ha\tbr{h}a-ni	&	\cg{(luma <Kay.)}	&	\cg{(kuki-ni)}	&	{\hr}bi\tbr{h}u	&	{\hr}haa\\ \midrule
%\tsc{Ambon-Seram} *R > **l\\ \midrule
Kayeli	&		&	{\hr}la\tbr{l}a	&	{\hr}\tbr{l}uma	&	{\hr}lo\tbr{l}i-n	&	{\hr}be\tbr{l}o	&	{\hr}laa-ni\\
Boano	&		&		&	{\hr}\tbr{n}uma	&	{\hr}lu\tbr{\und{l}}i	&	{\hr}he\tbr{n}u	&	{\hr}naa\tbr{n}\\
Luhu	&		&	{\hr}la\tbr{l}a-i	&	{\hr}\tbr{l}uma	&	{\hr}lu\tbr{l}i	&		&	{\hr}laʔa-ne\\
Piru	&	{\hr}ha\tbr{l}at	&	{\hr}la\tbr{l}a	&	{\hr}\tbr{l}uma	&	{\hr}lu\tbr{l}i \it{(+t)}	&	{\hr}he\tbr{l}u	&	{\hr}laia\tbr{l}-ai-n\su{†}\\
Batumerah	&		&	{\hr}la\tbr{l}a-i	&		&		&		&	\\
Wakasihu	&		&		&		&		&		&	{\hr}lai\tbr{l}-u\\
Larike	&	{\hr}ha\tbr{l}at-a	&	{\hr}la\tbr{l}a-u	&	{\hr}\tbr{\und{d}}uma	&	{\hr}du\tbr{\und{d}}i \it{(+t)}	&	{\hr}hi\tbr{\und{d}}o	&	{\hr}lae\tbr{\und{d}}-u\\
Asilulu	&	{\hr}ha\tbr{l}at	&	{\hr}la\tbr{l}a	&	{\hr}\tbr{l}uma	&	{\hr}lu\tbr{l}i- \it{(+t)}	&	{\hr}he\tbr{l}u	&	{\hr}laa\tbr{l}\\
Kaitetu	&		&	{\hr}la\tbr{l}a	&	{\hr}\tbr{l}uma	&	{\hr}lu\tbr{l}i	&		&	{\hr}laa\tbr{l}\\
Paulohi	&	{\hr}ha\tbr{l}at-e	&	{\hr}la\tbr{l}a	&	{\hr}\tbr{\und{ʀ}}uma	&	{\hr}ʀi\tbr{\und{ʀ}}i-a	&	{\hr}he\tbr{\und{ʀ}}u-i	&	{\hr}laa\tbr{l}-e\\
Kaibobo	&	{\hr}ha\tbr{l}at	&	{\hr}la\tbr{l}a	&	{\hr}\tbr{\und{r}}uma	&	{\hr}ru\tbr{\und{r}}i \it{(+t)}	&	{\hr}hu\tbr{\und{r}}u	&	{\hr}laia\tbr{l}-ai-n\su{†}\\
Amahai	&	{\hr}ho\tbr{l}aʔ-olo	&	{\hr}la\tbr{l}a	&	{\hr}\tbr{\und{r}}uma	&	{\hr}ri\tbr{\und{r}}i	&	{\hr}ho\tbr{\und{r}}o	&	\\
Kamarian	&	{\hr}ha\tbr{l}at	&	{\hr}la\tbr{l}a	&	{\hr}\tbr{\und{r}}uma	&	{\hr}ru\tbr{\und{r}}i \it{(+t)}	&	{\hr}he\tbr{\und{r}}u	&	{\hr}laia\tbr{r}-ai-n\su{†}\\
Haruku	&	{\hr}ha\tbr{l}at-o	&	{\hr}la\tbr{l}a	&	{\hr}\tbr{\und{r}}uma	&	{\hr}ru\tbr{\und{r}}i	&	{\hr}ho\tbr{\und{r}}u-ele	&	{\hr}lai\tbr{l}-o\\
Saparua	&	{\hr}ha\tbr{l}at-ol	&	{\hr}lala-l	&	{\hr}\tbr{l}uma	&	{\hr}li\tbr{l}i-n	&	{\hr}hu\tbr{l}u	&	{\hr}laia\tbr{l}-ai-n\su{†}\\
Nusalaut	&	{\hr}ha\tbr{l}at-ol	&	{\hr}lala-l	&	{\hr}\tbr{l}uma	&	{\hr}ri\tbr{\und{r}}i-ɲo	&	{\hr}hu\tbr{\und{r}}u	&	{\hr}laia\tbr{l}-ai-n\su{†}\\
N. Wemale	&	{\hr}ha\tbr{l}at-e	&	{\hr}la\tbr{l}a	&	{\hr}\tbr{l}uma	&	{\hr}lu\tbr{l}i	&	{\hr}pe\tbr{l}u-i	&	{\hr}leya\tbr{l}-e\\
Sepa	&	{\hr}ha\tbr{r}at-o	&	{\hr}la\tbr{y}a	&	{\hr}\tbr{r}uma	&	\cg{(toi)}	&	{\hr}ho\tbr{r}u-o	&	{\hr}lia\tbr{l}-o\\
Tamilou	&	{\hr}ha\tbr{r}at-o	&	{\hr}la\tbr{y}a	&	{\hr}\tbr{y}uma	&	\cg{(toi)}	&	{\hr}ho\tbr{y}u-o	&	{\hr}lea\tbr{l}-o\\
Tehoru	&	{\hr}ha\tbr{r}at-o	&	{\hr}la\tbr{y}a	&	{\hr}uma	&	\cg{(toi)}	&	{\hr}hoiʔ-o	&	\cg{(layar-o <Mly)}\\
N. Alune	&	{\hr}ba\tbr{l}at-e	&	{\hr}la\tbr{l}akwe	&	{\hr}\tbr{l}uma	&	{\hr}lu\tbr{l}i \it{(+t)}	&	{\hr}be\tbr{l}u-kwe	&	{\hr}lae\tbr{l}-e\\
W. Alune	&		&		&	{\hr}\tbr{r}uma	&	{\hr}du\tbr{\und{d}}i \it{(+t)}	&	{\hr}be\tbr{\und{d}}u-ke	&	\\
Yalahatan	&	{\hr}ha\tbr{l}at-e	&	{\hr}ra\tbr{\und{r}}a-ne	&	{\hr}\tbr{l}uma	&	{\hr}ru\tbr{\und{r}}i-n	&	{\hr}he\tbr{l}u {\tl} fe\tbr{l}u	&	{\hr}leʔa\tbr{l}-e\\
S. Nuaulu	&	{\hr}hanat-e	&	\cg{(nasi)}	&	{\hr}\tbr{n}uma	&	{\hr}u\tbr{n}i-e	&	{\hr}ho\tbr{n}u-e	&	{\hr}na\tbr{n}-e\\
Manusela	&	\cg{(lia sopu)}	&	\cg{(lasi-a)}	&	{\hr}\tbr{l}uma	&	{\hr}hu\tbr{l}i-kata	&	{\hr}ho\tbr{l}u-a	&	\\
Liana-Seti	&	{\hr}fa\tbr{r}at	&	{\hr}la\tbr{h}a	&	\cg{(ufaiya)}	&	{\hr}en-lui\tbr{l}-a	&	{\hr}fo\tbr{l}a	&	{\hr}laiya\tbr{l}-a\\ \midrule
Banda	&	{\hr}fa\tbr{r}at	&	{\hr}ra\tbr{r}a	&	{\hr}\tbr{r}umo	&	{\hr}ri\tbr{r}i-	&	{\hr}fe{\tl}fe\tbr{r}uŋ	&	{\hr}lea\tbr{r}\\ \midrule
%\tsc{Seram-Tanimbar-Bomberai} *R > **r\\ \midrule
Bobot	&		&	{\hr}-la\tbr{w}a	&	{\hr}uma	&	{\hr}-luin	&	{\hr}vo{\tl}voi	&	\\
Geser	&	{\hr}wa\tbr{r}at	&	{\hr}ra\tbr{r}a	&	{\hr}\tbr{r}uma	&	{\hr}ru\tbr{r}i	&	{\hr}wou{\tl}wou	&	{\hr}la\tbr{r}\\
Watubela	&		&	{\hr}la\tbr{l}ak	&	{\hr}\tbr{l}umak	&	{\hr}lu\tbr{l}i	&	{\hr}oku{\tl}oku	&	\\
Kowiai	&		&	{\hr}ra\tbr{r}a	&	\cg{(san)}	&	{\hr}ru\tbr{r}	&	{\hr}ɸowa	&	\cg{(ambaʔ)}\\
Yamdena	&	{\hr}ba\tbr{r}at	&	{\hr}da\tbr{r}-e	&	\cg{(das-e)}	&	{\hr}du\tbr{r}i \it{(+t)}	&	{\hr}be{\tl}be\tbr{r}y	&	{\hr}la\tbr{r}-e\\
Kei	&	{\hr}va\tbr{r}at	&	{\hr}la\tbr{r}	&	\cg{(raha)}	&	{\hr}lu\tbr{r}i-n	&	{\hr}vi\tbr{r}un	&	{\hr}laa\tbr{r}\\
Uruangnirin	&		&	{\hr}na\tbr{r}ah	&	\cg{(kapala)}	&	{\hr}nu\tbr{r}i	&	{\hr}per{\tl}pe\tbr{r}i	&	\cg{(kabak)}\\
Sekar	&		&	{\hr}ra\tbr{r}a	&	{\hr}\tbr{r}uma	&	{\hr}ru\tbr{r}i-r \it{(+t)}	&	{\hr}be\tbr{r}i	&	{\hr}ra\tbr{r}\\
Arguni	&		&	{\hr}rá\tbr{r}e	&	{\hr}\tbr{r}úme	&	{\hr}ru\tbr{r}	&	{\hr}pu{\tl}pua\tbr{r}u	&	{\hr}re\tbr{r}\\
Irarutu	&		&	\cg{(wams)}	&	{\hr}\tbr{r}um	&	{\hr}ru\tbr{r}	&	\cg{(bunat)}	&	\cg{(maru)}\\
		\lspbottomrule
	\end{tabular}
		\begin{tablenotes}\tnsize
			\item[†] {Piru, Kaibobo, Kamarian, Saparua, Nusalaut ‘mast of a boat’; literally ‘wood for sail’.}
		\end{tablenotes}
	\end{threeparttable}
\end{adjustbox}
\end{table}

\begin{table}[p]\centering\tcp{Split of *j distinguishing languages of central Maluku}{07-12}
\begin{adjustbox}{totalheight=\textheight}
	\begin{threeparttable} %\st{2pt}
	\begin{tabular}{lll@{ }lll}\lsptoprule
*gloss	&	name	&	\mc{2}{l}{gall(bladder)}	&	younger sibling\su{‡}	&	nose\\
PMP	&	*ŋa\tbr{j}an	&	\mc{2}{l}{*qapə\tbr{j}u}	&	*hua\tbr{j}i	&	*(ŋ)i\tbr{j}uŋ\\ \midrule
%\tsc{Sula-Buru}\\ \midrule
Mangole	&	{\hr}ŋaa	&	{\hr}n-peu	&		&	\cg{(fuki)}	&	\cg{(nʧumi)}\\
Sanana	&	{\hr}naa	&	{\hr}m-peu	&		&	\cg{(fuk)}	&	{\hr}neʔi\\
Lisela	&	{\hr}ŋaa-n	&	{\hr}peu-n	&		&	{\hr}uwai	&	{\hr}ŋee-n\\
Buru	&	{\hr}ŋaa-n	&	{\hr}peu-n	&		&	{\hr}wai \cg{(ySib SS)}	&	{\hr}ŋee-n\\
Ambelau	&	{\hr}nea-mo	&	{\hr}feo-ni	&		&	{\hr}βai-o, eβai-no	&	{\hr}neini\\ \midrule
%\tsc{Ambon-Seram}\\ \midrule
Kayeli	&	{\hr}naa-ni (< SB)	&		&		&	{\hr}wai-n (< SB)	&	{\hr}neeɲi (< SB)\\
Boano	&	{\hr}na\tbr{n}a	&		&		&	{\hr}a\tbr{n}i	&	{\hr}ni\tbr{n}u\\
Piru	&	{\hr}na\tbr{l}a	&	{\hr}o\tbr{l}u 	&	‘liver’	&	{\hr}wa\tbr{l}i	&	{\hr}ni\tbr{n}u\\
Larike	&	{\hr}na\tbr{l}a	&	{\hr}í\tbr{\und{d}}ona	&		&	{\hr}wa\tbr{l}e	&	\cg{(disu)}\\
Asilulu	&	{\hr}na\tbr{l}a-	&	{\hr}we\tbr{l}u	&		&	{\hr}wa\tbr{l}i-	&	{\hr}ni\tbr{n}u\\
Kaitetu	&	{\hr}na\tbr{l}a	&	{\hr}we\tbr{l}ui 	&	‘liver’	&	{\hr}wa\tbr{l}i	&	\cg{(lisu)}\\
Paulohi	&	{\hr}na\tbr{l}a	&		&		&	{\hr}wa\tbr{\und{ʀ}}i	&	{\hr}i\tbr{\und{ʀ}}i\\
Kaibobo	&	{\hr}na\tbr{l}a	&	{\hr}o\tbr{\und{r}}u	&	‘liver’	&	{\hr}a\tbr{\und{r}}i	&	{\hr}i\tbr{\und{r}}u\\
Kamarian	&	{\hr}na\tbr{l}a	&	{\hr}ho\tbr{\und{r}}ui 	&	‘liver’	&	{\hr}wa\tbr{\und{r}}i	&	{\hr}hi\tbr{\und{r}}u\\
Haruku	&	{\hr}na\tbr{l}a-i	&	{\hr}wo\tbr{\und{r}}ui	&	‘liver’	&	{\hr}wa\tbr{\und{r}}i	&	{\hr}i\tbr{\und{r}}u\\
Saparua	&	{\hr}na\tbr{l}a	&	{\hr}wo\tbr{\und{r}}u-lo 	&	‘liver’	&	{\hr}wa\tbr{\und{r}}i	&	{\hr}i\tbr{l}i\\
Nusalaut	&	{\hr}na\tbr{l}a	&	{\hr}wo\tbr{\und{r}}uil	&	‘liver’	&	{\hr}wa\tbr{\und{r}}i	&	{\hr}i\tbr{\und{r}}i\\
Sepa	&	{\hr}ni\tbr{y}a-m\su{†}	&	{\hr}hu\tbr{\und{r}}u-n	&		&	{\hr}wa\tbr{\und{r}}i	&	{\hr}i\tbr{\und{r}}i-n\\
Tamilou	&	{\hr}ni\tbr{y}a-m\su{†}	&	{\hr}ho\tbr{\und{l}}u-n	&		&	{\hr}wa\tbr{\und{l}}i-n	&	{\hr}u\tbr{\und{l}}i\\
Tehoru	&	{\hr}ni\tbr{y}a-m\su{†}	&	{\hr}f\und{oi}-n\su{†}	&		&	{\hr}wa\tbr{\und{l}}i-n	&	\\
N. Wemale	&	{\hr}na\tbr{l}a	&	{\hr}ho\tbr{l}u	&		&	{\hr}wa\tbr{l}i	&	{\hr}i\tbr{l}i\\
S. Wemale	&	{\hr}naa (*l > {\0} /a{\gap}a)	&		&	{\hr}wa\tbr{l}e	&	\\
N. Alune	&	{\hr}na\tbr{l}a-ne	&	{\hr}so\tbr{l}u	&		&	{\hr}kwa\tbr{l}i muli	&	{\hr}i\tbr{n}u\\
Yalahatan	&	{\hr}na\tbr{l}a-i	&		&		&	{\hr}wa\tbr{l}i-n	&	\cg{(nua)}\\
Atamano	&	{\hr}na\tbr{l}a	&	{\hr}o\tbr{l}u	&		&		&	\\
S. Nuaulu	&	{\hr}na\tbr{n}a-i	&	\cg{(moni-ni)}	&	{\hr}wa\tbr{n}i-ni	&	\cg{(oko-e)}\\
Manusela	&	{\hr}la\tbr{l}a-mu	&	{\hr}wo\tbr{l}u-ni	&		&		&	\cg{(nita)}\\
Liana-Seti	&	{\hr}na\tbr{l}a-k	&	{\hr}hu\tbr{l}aina	&		&	\cg{(foto-nin)}	&	{\hr}ni\tbr{n}u-\\ \midrule
Banda	&	{\hr}na\tbr{l}a-n	&		&		&	{\hr}wa\tbr{l}i-n	&	{\hr}nu\tbr{l}un\\ \midrule
%\tsc{Seram-Tanimbar-Bomberai}\\ \midrule
Bobot	&	{\hr}-na\tbr{l}an	&	{\hr}fo\tbr{l}i	&		&	{\hr}wa\tbr{l}i-n	&	{\hr}-i\tbr{l}in\\
Masiwang	&	{\hr}na\tbr{l}an	&	{\hr}ho\tbr{l}i	&		&	{\hr}wa\tbr{l}i	&	{\hr}ni\tbr{n}in\\
Geser	&	{\hr}ŋa\tbr{s}an	&	{\hr}fo\tbr{l}i	&		&	{\hr}a\tbr{l}i	&	{\hr}\tbr{s}o\\
Gorom	&	{\hr}ŋa\tbr{s}an	&		&		&	{\hr}a\tbr{l}i	&	{\hr}i\tbr{s}o, \tbr{s}ue\\
Watubela	&	{\hr}ŋa\tbr{h}an	&	{\hr}fe\tbr{l}i-n	&		&		&	{\hr}\tbr{h}u-)\\
Kowiai	&	{\hr}ne\tbr{s}a	&		&		&	\cg{(ɸuta)}	&	\cg{(siʔai)}\\
Kur	&	{\hr}ŋaa-m	&		&		&		&	\\
Yamdena	&	{\hr}ŋa\tbr{r}-e	&	{\hr}fe\tbr{r}u-	&		&	{\hr}wai-ny	&	{\hr}i\tbr{r}i-ŋw\\
Fordata	&	{\hr}na\tbr{r}a-n	&	{\hr}fe\tbr{l}ur-n	&		&	{\hr}wa\tbr{r}i-n	&	{\hr}ni\tbr{r}u-n\\
Uruangnirin	&	{\hr}ga\tbr{r}an, ŋga\tbr{r}an	&	\cg{(ŋgeram)}	&		&	{\hr}wa\tbr{r}i-n	&	{\hr}ni\tbr{r}iŋgain\\
Onin	&	{\hr}ga\tbr{r}a-r	&	{\hr}ma\tbr{r}ik	&	‘bitter’	&	{\hr}wa\tbr{r}i-r	&	{\hr}i\tbr{r}i-r\\
Sekar	&	{\hr}ga\tbr{r}a-r	&	{\hr}ma\tbr{r}ik	&	‘bitter’	&		&	\cg{(idum < Mly.)}\\
Arguni	&	{\hr}gá\tbr{d}in	&	{\hr}ma\tbr{r}i	&	‘bitter’	&	{\hr}twá\tbr{r}i	&	\cg{(sigyó)}\\
Irarutu	&	\cg{(no)}	&	{\hr}-fi\tbr{r}-	&	‘bile’	&	\cg{(nefut)}	&	{\hr}wegu\tbr{r}\\
		\lspbottomrule
	\end{tabular}
		\begin{tablenotes}\tnsize
			\item[†]	{See \srf{10-01-02} for discussion of \tsc{East Littoral} words which
								show **l > \it{y}.}
			\item[‡]	{Local glosses include ‘sibling,
								younger sibling, younger sibling (same sex),
								brother-in-law (mostly ZH, sometimes WB)’.}
		\end{tablenotes}
	\end{threeparttable}
\end{adjustbox}
\end{table}
